\chapter{The Modal Mindset}\label{ch:vim-modal}

\section{Why Modes}\label{sec:vim-why-modal}

Every other editor you have used is \vocab{modeless}: the keyboard always
inserts text, and commands are reached by holding modifiers
(\key{<C-c>}, \key{<C-v>}, \key{<C-S-k>}) or by leaving the keyboard for the
mouse. This is a reasonable design, but it wastes the most valuable real estate
on the machine --- the twenty-six letter keys --- on a single operation.

\begin{intuition}
	\vocab{Modal editing} observes that you spend far more time \emph{reading and
		restructuring} code than typing it fresh. So Vim gives the letter keys a
	second life: in \vocab{Normal mode} they are commands, and you enter
	\vocab{Insert mode} only for the short bursts when you are genuinely producing
	new characters.
\end{intuition}

The consequence is that commands become \emph{composable} in a way that
modifier-key shortcuts never can be. \key{<C-S-k>} deletes a line and that is
all it will ever do; \key{d} is an operator that combines with any motion, so
\key{d} \key{w}, \key{d} \key{\$}, \key{d} \key{i} \key{(} and \key{d} \key{/}
\key{f} \key{o} \key{o} \key{<CR>} all follow from one key. We will make this
precise in \autoref{ch:vim-grammar} --- it is the single idea worth taking away from
this note.

\begin{moral}
	Normal mode is home. You should be in Insert mode only while text is actually
	appearing on screen, and you should leave it the moment you stop.
\end{moral}

\section{The Modes}\label{sec:vim-the-modes}

\begin{definition}[Normal mode]
	The default mode, entered with \key{<Esc>}. Keys are \vocab{commands}: they
	move the cursor, delete, copy, paste and repeat. If you are ever unsure what
	mode you are in, press \key{<Esc>} twice.
\end{definition}

\begin{definition}[Insert mode]
	Typed characters are inserted into the buffer, as in any ordinary editor.
	Entered with \key{i}, \key{a}, \key{o} and friends (\autoref{sec:vim-insert}),
	left with \key{<Esc>} or \key{<C-[>}.
\end{definition}

\begin{definition}[Visual mode]
	Select a region first, then act on it --- the reverse of Normal mode's
	verb-then-object order. Three flavours:
	\begin{keys}[0.16]
		\key{v}      & \vocab{charwise}: select character by character                 \\
		\key{V}      & \vocab{linewise}: select whole lines                            \\
		\key{<C-v>}  & \vocab{blockwise}: select a rectangle, for column edits         \\
	\end{keys}
\end{definition}

\begin{definition}[Command-line mode]
	Entered with \key{:} (Ex commands such as \cmd{:write}), \key{/} or \key{?}
	(search), or \key{!} (filter through a shell command). One line is typed at the
	bottom of the screen and executed with \key{<CR>}; \key{<Esc>} aborts.
\end{definition}

\begin{definition}[Replace mode]
	Entered with \key{R}. Typed characters overwrite existing ones instead of
	pushing them right. Rarely needed; \key{r} (replace exactly one character) is
	the version you will actually use.
\end{definition}

\begin{definition}[Terminal mode]
	A Neovim addition: inside a \cmd{:terminal} buffer, keystrokes go to the shell.
	Return to Normal mode with \key{<C-\textbackslash>} \key{<C-n>}. See
	\autoref{sec:nvim-terminal}.
\end{definition}

\begin{remark}
	Vim also has \vocab{Operator-pending mode} --- the sliver of time after you
	press \key{d} but before you have said \emph{what} to delete. You will never
	enter it deliberately, but it is why \key{d} \key{d} works: the second \key{d}
	is parsed as ``the whole line'' only because Vim was waiting for a motion.
\end{remark}

\hrulebar

\section{Survival Kit: Open, Edit, Save, Quit}\label{sec:vim-survival}

Enough to open a file, change it, and get out alive.

\begin{keys}[0.22]
	\cmd{nvim file.tex}  & Open a file from the shell                                   \\
	\key{i}              & Enter Insert mode; type something                            \\
	\key{<Esc>}          & Back to Normal mode                                          \\
	\cmd{:w<CR>}         & \vocab{Write} (save) the file                                \\
	\cmd{:q<CR>}         & \vocab{Quit}                                                 \\
	\cmd{:wq<CR>}        & Write and quit (\key{Z} \key{Z} is the Normal-mode shorthand) \\
	\cmd{:q!<CR>}        & Quit, discarding all changes                                 \\
	\cmd{:w newname<CR>} & Write to a different file                                     \\
	\key{u}              & Undo                                                         \\
	\key{<C-r>}          & Redo                                                         \\
\end{keys}

\begin{note}[The \texttt{!} suffix]
	Appending \key{!} to an Ex command means \emph{force}, or ``yes, I am sure''.
	\cmd{:q} refuses to quit a modified buffer; \cmd{:q!} throws the changes away.
	\cmd{:qa!} does it for every open buffer at once.
\end{note}

\begin{note}[If you are stuck]
	Three situations trap beginners, and all three have the same escape:
	\begin{enumerate}[(i)]
		\item You pressed \key{<C-s>} out of habit and the terminal froze --- press
		      \key{<C-q>} to resume.
		\item You are in the middle of a half-typed command --- \key{<Esc>}.
		\item You entered Ex mode (\key{Q} in Vim) and see a \texttt{:} prompt that will
		      not go away --- type \cmd{:visual<CR>}. Neovim removed this trap entirely.
	\end{enumerate}
\end{note}

\begin{tryit}
	Open a scratch file with \cmd{nvim /tmp/scratch.txt}, type a few lines, save
	with \cmd{:w}, quit with \cmd{:q}, reopen it, and delete everything you wrote
	with \key{u}. Do this until saving and quitting requires no thought at all ---
	the rest of the note assumes it.
\end{tryit}

\section{Reading the Manual}\label{sec:vim-help}

Vim's documentation is unusually good, and it is available offline, inside the
editor, with no network round trip. Learning to navigate it is worth more than
any cheat sheet.

\begin{keys}[0.30]
	\cmd{:help}                & The table of contents                                        \\
	\cmd{:help dd}             & Help for the Normal-mode command \key{d} \key{d}             \\
	\cmd{:help i\_CTRL-r}      & Help for \key{<C-r>} \emph{in Insert mode}                   \\
	\cmd{:help v\_o}           & Help for \key{o} \emph{in Visual mode}                       \\
	\cmd{:help 'number'}       & Help for an \emph{option} --- note the single quotes          \\
	\cmd{:help :write}         & Help for an Ex command --- note the leading colon             \\
	\cmd{:helpgrep <pattern>}  & Full-text search across all help files                       \\
	\key{<C-]>}                & Follow the tag under the cursor                              \\
	\key{<C-o>}                & Jump back where you came from                                \\
\end{keys}

So the mode you are asking about is encoded in the prefix: \cmd{i\_} for Insert,
\cmd{v\_} for Visual, \cmd{c\_} for the command line, \cmd{o\_} for
operator-pending, and no prefix for Normal. Tab completion works on the help
argument, which is the fastest way to discover the prefixes.

\begin{eg}
	You have forgotten what \key{g} \key{U} does. Type \cmd{:help gU} and read.
	You want to know why your search is case-insensitive: \cmd{:help 'ignorecase'}.
	You saw \key{<C-a>} in someone's config and want to know if it is safe to
	remap: \cmd{:help CTRL-A}.
\end{eg}

\begin{note}[\texttt{vimtutor}]
	Vim ships with a thirty-minute interactive tutorial. Run \cmd{vimtutor} in a
	shell (or \cmd{:Tutor} inside Neovim). It is the single highest-return half
	hour available to a beginner, and this chapter is not a substitute for it.
\end{note}

\section{Options}\label{sec:vim-options}

Behaviour is controlled by \vocab{options}, set with \cmd{:set}. Three forms
matter:

\begin{keys}[0.30]
	\cmd{:set number}        & Turn a boolean option on                    \\
	\cmd{:set nonumber}      & Turn it off (prefix \cmd{no})               \\
	\cmd{:set number!}       & Toggle it (suffix \cmd{!})                  \\
	\cmd{:set number?}       & Ask its current value (suffix \cmd{?})      \\
	\cmd{:set tabstop=4}     & Set a numeric or string option              \\
	\cmd{:setlocal ...}      & Apply to the current buffer/window only     \\
\end{keys}

A handful you will want on day one:

\begin{keys}[0.30]
	\cmd{'number'}, \cmd{'relativenumber'} & Absolute and relative line numbers; together they give the \vocab{hybrid} display that makes counts like \key{7} \key{j} easy to read off \\
	\cmd{'ignorecase'}, \cmd{'smartcase'}  & Case-insensitive search, unless the pattern contains a capital                                                                          \\
	\cmd{'expandtab'}, \cmd{'shiftwidth'}  & Insert spaces instead of tabs; how far \key{>} \key{>} indents                                                                          \\
	\cmd{'wrap'}, \cmd{'linebreak'}        & Soft-wrap long lines, and wrap at word boundaries                                                                                      \\
	\cmd{'scrolloff'}                      & Keep $n$ lines of context above and below the cursor                                                                                   \\
	\cmd{'clipboard'}                      & Set to \cmd{unnamedplus} to share yanks with the system clipboard (\autoref{sec:vim-registers})                                            \\
\end{keys}

Typed at the \key{:} prompt, an option lasts for the session only. To make it
permanent it goes in your configuration file --- \file{\textasciitilde/.vimrc}
for Vim, \file{\textasciitilde/.config/nvim/init.lua} for Neovim, where the Lua
spelling is \cmd{vim.opt.number = true} (\autoref{sec:nvim-lua}).

\begin{remark}
	Neovim already switches on many options that Vim leaves off for compatibility
	with 1976: \cmd{'nocompatible'}, \cmd{'incsearch'}, \cmd{'hlsearch'},
	\cmd{'autoindent'}, syntax highlighting and filetype plugins are all default.
	A Neovim configuration is therefore much shorter than the equivalent
	\file{.vimrc}; see \cmd{:help nvim-defaults} for the full list.
\end{remark}
